% Generated draft: requires mathematical and source review.
\documentclass[11pt]{article}
\usepackage[margin=25mm]{geometry}
\usepackage{fontspec}
\setmainfont{Latin Modern Roman}
\usepackage{amsmath,amssymb,mathtools,unicode-math}
\setmathfont{Latin Modern Math}
\usepackage{enumitem,needspace,xurl,hyperref}
\setlength{\parindent}{0pt}
\setlength{\parskip}{5pt}
\setlength{\emergencystretch}{3em}
\newcommand{\N}{\mathbb{N}}
\newcommand{\Z}{\mathbb{Z}}
\newcommand{\Q}{\mathbb{Q}}
\newcommand{\R}{\mathbb{R}}
\newcommand{\C}{\mathbb{C}}
\newcommand{\F}{\mathbb{F}}
\newcommand{\A}{\mathbb{A}}
\newcommand{\PP}{\mathbb{P}}
\newcommand{\E}{\mathbb{E}}
\newcommand{\eps}{\varepsilon}
\newcommand{\dd}{\,\mathrm{d}}
\begin{document}
\section*{UnsolvedMath ID: 30002184}
\textbf{Source problem code:} catalogue entry
\subsection{Definitions and mathematical statement}
The supplied question is: ``How long may it take for a crossing event to disappear in dynamical percolation?'' The record does not specify enough model data to determine a unique persistence-time problem. The following gives a conditional formulation, with its unspecified parameters displayed explicitly.

Fix a graph \(G\), a subgraph \(D\), and two designated boundary sets \(B_1,B_2\subseteq V(D)\). Specify either site or bond percolation. Let \((\omega_t)_{t\geq 0}\) be the chosen dynamical-percolation process, with its percolation parameter, update mechanism, update rates, and boundary conditions specified. Define the crossing event
\[
 A=\{\omega:\text{there is an open path in }D\text{ joining }B_1\text{ to }B_2\}.
\]
Here openness is interpreted according to the selected site or bond model. For a specified initial distribution \(\mu\) satisfying \(\mu(A)=1\), define the first disappearance time and its tail by
\[
 T_A=\inf\{t\geq 0:\omega_t\notin A\},\qquad
 S_{A,\mu}(t)=\mathbb P_\mu(T_A>t),\quad t\geq 0,
\]
with \(\inf\varnothing=+\infty\). Conditioning a stationary initial distribution on \(A\) is one possible initial convention, but the source does not specify that convention.

\textbf{Conditional mathematical target.} For the dynamical-percolation model and crossing geometry intended by the source, determine the distribution of \(T_A\), or quantitative bounds and asymptotic behavior for \(S_{A,\mu}(t)\). If a family of crossing regions is intended, its indexing, time normalization, and order of limits must also be specified. This is a formulation template, not a claim that one bound should hold for every graph, parameter, geometry, update rule, or initial law. The source does not identify whether the requested result is an expectation, a tail estimate, a sharp asymptotic law, or a statement uniform over scales.

The supplied background discusses critical dynamical percolation and box crossings, but the question itself does not impose criticality or identify a lattice, dimension, or box. Its dated literature assessment reports superpolynomial decay for box-crossing persistence and exponential tails for the first exceptional infinite-cluster time. The latter is a different stopping-time question and is not identified here with \(T_A\). The record labels the problem partially solved; this label and the reported bounds are supplied curation, not independently verified mathematical conclusions.
\subsection{Short English statement}
Starting with an open crossing, ask how long it remains present under dynamical percolation. A precise persistence-tail question requires the graph, crossing geometry, percolation parameter, dynamics, initial law, and time scale. These are not fixed by the supplied record.
\subsection{Sources}
\begin{itemize}
\item Local UnsolvedMath registry, permanent ID \(30002184\), source problem code \texttt{catalogue entry}, title ``Persistence Times of Dynamical-Percolation Crossings.'' The supplied record belongs to the Oberwolfach Reports Open Problems collection and gives proposed year \(2012\); no individual proposer is identified.
\item Supplied report citation: \emph{Scaling Limits in Models of Statistical Mechanics} (2013), Oberwolfach Reports, DOI \texttt{10.4171/owr/2012/44}, pp.\ 2657--2745. \url{https://doi.org/10.4171/owr/2012/44}. The background describes Oberwolfach Reports as documenting mathematical workshops and their open-problem discussions. The supplied citation uses 2013, while the DOI and proposed-year metadata refer to 2012; these details are retained without reconciliation.
\item The record labels its statement \texttt{exact} and reports that its mathematical content was checked against the official Oberwolfach report, with notation, context, and question form normalized without changing the source-level claim. This is a supplied verification assertion; no independent source inspection is claimed here. No particular problem page within the cited report is supplied.
\item Supplied literature reference: Hammond--Mossel--Pete, ``exit-time tails.'' \url{https://doi.org/10.1214/EJP.v17-2229}.
\item Supplied literature reference: Hammond--Pete--Schramm, ``exceptional-set local time.'' \url{https://arxiv.org/abs/1208.3826}.
\item Literature assessment checked at \textbf{2026-08-21}: the supplied status is \texttt{partially\_solved}. The assessment reports superpolynomial decay for persistence of a box crossing in critical dynamical percolation and exponential tails for the first exceptional infinite-cluster time, while stating that these results do not determine a sharp general tail law. This is supplied dated curation derived from the OWR collection, not independently verified. The record expressly requests independent verification of time-sensitive status and cited claims.
\end{itemize}
\end{document}
